\subsection*{Note (4)}
The fourth note, on pp.~54--55, expands the motivic direction attached to
\S3, p.~17. Grothendieck describes the arithmetic structure on
$S\ell(2,\mathbb Z)$ as a ``machine à motifs'', recalls experiments begun
in 1981, and specifies the input homomorphisms into $G_o(\mathbb Z)$ and
the group-scheme extensions used for the resulting Galois
representations. The tautological case leads in his account to an
extension of $G\ell(2)_{\mathbb Z}$ by a two-dimensional torus and to the
representations associated with the two named complex multiplication
fields. He records that this construction principle has not yet been
followed through. The passage from that unfinished exploration to the
five-part priority list on p.~55 makes the program's current order
explicit while calling the order provisional: construction of the
Teichmüller tower, its profinite automorphism group and the Galois
subgroup, the motivic machine, the anabelian dictionary and fundamental
conjecture, and Fermat's problem. Thus items a) and b) return to \S2,
item c) to the Jacobian representations of \S3, and item d) to the broader
reconstruction of note (3). The parenthesis saying that the fundamental
conjecture may be less ``hors de portée'' than supposed continues the
explicit change of mind recorded in note (1). The ``paradigme profini'',
inertia subgroups, positive characteristic, and the absolute ring are
named here as parts of the intended expansion, without a proposed full
formulation being supplied.
(\href{https://grothendieck.umontpellier.fr/web/viewer.html?file=../119.pdf\#page=57}{Copy 1, pp.~54--55, beginning at site 56};
\href{https://grothendieck.umontpellier.fr/web/viewer.html?file=../119.pdf\#page=58}{priority list on p.~55, site 57};
\href{https://grothendieck.umontpellier.fr/web/viewer.html?file=../119.pdf\#page=20}{call in \S3, p.~17, site 19}.)
